\section*{About this document}
\markboth{About this document}{About this document}

\textbf{Assumptions made for this document.} Study area: Semarang--Demak. Team: four students with one
radar specialist. Time: six weeks. Change any of these and the method stays; only
Sections~\ref{sec:data} and~\ref{sec:feas} need updating.

\subsection*{How to read this document}

Every item that matters has a tag: a letter for its kind and a circled number. Tags let you follow
one thread through the whole document. For example, \Q{2} (which sources?) is answered by stage \M{6},
checked by \V{3}, built in component \Sx{4}, and delivered as output \Ox{3}. Section~\ref{sec:trace}
lists every such link.

\begin{center}
\small\renewcommand{\arraystretch}{1.0}
\begin{tabular}{@{}llp{62mm}l@{}}
\toprule
\textbf{Tag} & \textbf{Kind} & \textbf{Meaning} & \textbf{Defined in}\\*
\midrule
\Q{1}--\Q{3} & Question & What the research must answer & Section 2\\
\D{1}--\D{9} & Data & A data set and where it comes from & Section~\ref{sec:data}\\
\M{1}--\M{9} & Method & One stage of the pipeline & Section~\ref{sec:method}\\
\V{1}--\V{7} & Validation & A check with a pass mark & Section~\ref{sec:rigour}\\
\R{1}--\R{8} & Risk & What could go wrong, and the response & Section~\ref{sec:feas}\\
\U{1}--\U{6} & User need & What a user must be able to do & Section~\ref{sec:product}\\
\Sx{1}--\Sx{9} & Software & A component of the system & Section~\ref{sec:arch}\\
\Ox{1}--\Ox{5} & Output & Something the project delivers & Section 2\\
\bottomrule
\end{tabular}
\end{center}

Three kinds of boxes recur. \textcolor{accent}{\textsf{\textbf{In short}}} gives the point of a
section. \textcolor{green}{\textsf{\textbf{In plain words}}} explains an idea without jargon.
\textcolor{warm}{\textsf{\textbf{Honest limit}}} says what a result does not show.

\enlargethispage{2\baselineskip}
{\sffamily\footnotesize\hypersetup{linkcolor=ink}\setcounter{tocdepth}{1}\setlength{\parskip}{0pt}\makeatletter\renewcommand*\l@section{\@dottedtocline{1}{0em}{2.2em}}\makeatother\tableofcontents}
\clearpage
